\chapter{Alcohol Use Screening}

\section*{What Is This Test?}
Alcohol use screening is a structured clinical assessment that identifies individuals whose alcohol consumption places them at risk for alcohol-related harm. Unlike carbohydrate-deficient transferrin (CDT) or blood alcohol concentration, which are laboratory measures, screening tools are validated questionnaires completed by the patient or administered by the clinician. Adults should be screened at least annually in primary care, and adolescents when risk factors are present. A positive screen prompts a brief intervention and, when indicated, referral for specialized treatment. Two complementary questionnaires are widely used in the United States: the AUDIT (Alcohol Use Disorders Identification Test, developed by the WHO) and the shorter AUDIT-C, plus single-item sedheavier-drinking screens for very brief encounters.

\section*{Alternative Names}
\begin{itemize}
  \item AUDIT
  \item AUDIT-C
  \item CAGE
  \item CAGE-AID
  \item Single-Item Alcohol Screening
  \item Alcohol Use Disorders Identification Test
  \item Alcohol Screening Questionnaire
  \item Brief Alcohol Screening
\end{itemize}
\section*{Principle}
Screening uses validated questionnaires that map drinking patterns, frequency, and consequences onto DSM-5 criteria for alcohol use disorder and NIAAA risk thresholds for hazardous drinking. Each item is scored; the summed score is compared against a cutoff. AUDIT's 10 items cover consumption (questions 1--3), dependence symptoms (question 4--6), and harmful consequences (questions 7--10) \cite{saunders1993development}. AUDIT-C uses only the three consumption items \cite{bush1998audit}. A single-item screen (``How many times in the past year have you had X or more drinks in a day?'') uses sex-specific thresholds (5 for men, 4 for women) and a single occurrence as positive.

\section*{History}
AUDIT was developed by the World Health Organization in 1989 after a six-country validation study and published in 1993 \cite{saunders1993development}. It replaced the unstructured clinical history taking that missed most at-risk drinkers. The CAGE questionnaire (Cut down, Annoyed, Guilt, Eye-opener), introduced by Ewing in 1984, remains useful for detecting alcohol dependence but is less sensitive for at-risk drinking \cite{ewing1984detecting}. The U.S. Preventive Services Task Force has endorsed alcohol screening and brief behavioral counseling for all adults since 2004 \cite{uspstf2018alcohol}.

\section*{Why This Test Is Ordered}
\begin{itemize}
  \item Routine annual adult screening for risky drinking (USPSTF Grade B recommendation)
  \item Pre-operative or prenatal risk assessment
  \item Evaluation of patients with unexplained abnormal liver function tests, pancreatitis, or trauma
  \item Monitoring relapse risk in patients with known alcohol use disorder
  \item Adolescents and young adults when alcohol use is suspected
\end{itemize}

\section*{Primary vs Secondary Test}
Primary screening for unhealthy alcohol use in primary care and opportunistic clinical settings.

\section*{How This Test Is Performed}
The patient self-administers the questionnaire (paper or electronic) or answers the clinician. AUDIT, AUDIT-C require 2--5 minutes; single-item screens take less than a minute. No blood draw is taken, but clinicians often order companion laboratory tests (GGT, MCV, AST/ALT ratio, carbohydrate-deficient transferrin) when dependence is suspected. Adolescent screening uses S2BI or CRAFFT rather than adult tools.

\section*{What Preparation Is Needed}
None. The patient should be sober at the time of screening.

\section*{Contraindications and Precautions}
\begin{itemize}
  \item None. Confidentiality is essential: patients understate consumption if they fear legal or employment consequences. Results should be reviewed separately from the medical record if the patient prefers.
\end{itemize}
\section*{Risks}
None. Rare patients may feel embarrassed or judged.

\section*{What Happens After the Test}
A positive screen (AUDIT 8+ for men, 7+ for women; AUDIT-C 4+ for men, 3+ for women; CAGE 2+) prompts a brief negotiated interview, motivational feedback, and a management plan. Scores 20+ indicate probable dependence and warrant referral to addiction medicine.

\section*{Understanding Results}

\subsection*{Below Normal}
A score below the cutoff indicates low-risk or no alcohol use. The clinician should still ask about binge-drinking episodes not captured by score-based criteria, and re-screen at the next recommended interval (annually).

\subsection*{Above Normal}
\begin{itemize}
  \item AUDIT 8--15 (men) or 7--15 (women): Hazardous or harmful drinking; brief intervention is indicated
  \item AUDIT 16--19: Likely alcohol dependence; brief intervention plus monitoring
  \item AUDIT 20+: Definite alcohol dependence; referral to specialty addiction treatment
  \item AUDIT-C 4+ (men), 3+ (women): Positive screen, proceed to full AUDIT
  \item CAGE 2+: Probable alcohol use disorder
  \item Single-item: Any positive response warrants full AUDIT and brief intervention
\end{itemize}

\section*{Normal Results}
\begin{itemize}
  \item \textbf{AUDIT total:} 0--7 in men; 0--6 in women
  \item \textbf{AUDIT-C:} 0--3 in men; 0--2 in women
  \item \textbf{CAGE:} 0--1
  \item \textbf{Single-item screen:} Zero occasions in the past year
  \item \textbf{NIAAA low-risk drinking limits:} $\leq$ 4 drinks on any day and $\leq$ 14/week for men; $\leq$ 3 drinks on any day and $\leq$ 7/week for women
\end{itemize}

\section*{Follow-up Tests}
\begin{itemize}
  \item \textbf{Carbohydrate-deficient transferrin (CDT):} Validates heavy drinking over the prior 1--2 weeks
  \item \textbf{GGT and MCV:} Markers of chronic heavy use; GGT elevated in 70\% of heavy drinkers
  \item \textbf{AST/ALT ratio $>$ 2:} Suggests alcoholic hepatitis
  \item \textbf{Blood alcohol concentration:} Confirms acute intoxication
  \item \textbf{DSM-5 diagnostic interview:} Definitive diagnosis of alcohol use disorder
  \item \textbf{CIWA-Ar:} Monitors alcohol withdrawal severity if dependence is detected
\end{itemize}

\section*{References}
\bibliographystyle{unsrtnat}
\bibliography{references}

\clearpage
\section*{Regulatory Status and Authorization Date}
This chapter describes a test category, not a specific commercial product. FDA, CDSCO, EU IVDR, and MHRA approval or registration dates therefore cannot be assigned without the assay manufacturer, product name, instrument, and intended use. For laboratory-developed tests, an individual agency approval date may not exist. See the Preface for the interpretation of regulatory status.

\clearpage
